\subsection{Linear Second-Order Equations}

\begin{definition}{Second-Order ODE}{second-order-ode}
  A \textbf{second-order ODE} has highest derivative $d^2y/dx^2$.
  A \textbf{linear second-order ODE} has the form
  \[
    y''+p(x)y'+q(x)y=f(x).
  \]
  It is \textbf{homogeneous} if $f(x)=0$.
\end{definition}

\begin{example}[Classifying Second-Order Equations]
  The equation
  \[
    \frac{d^2y}{dx^2}+(\sin x)\frac{dy}{dx}+e^{3x^2}y=4\tan x
  \]
  is a nonhomogeneous linear second-order differential equation.
  The equation $y''+yy'=0$ is nonlinear and second-order.
\end{example}

\begin{theorem}{Existence and Uniqueness for Linear Second-Order Equations}{linear-second-order-existence-uniqueness}
  Consider the IVP
  \[
    \begin{cases}
      y''+p(x)y'+q(x)y=f(x), \\
      y(a)=b_0,\qquad y'(a)=b_1.
    \end{cases}
  \]
  If $p$, $q$, and $f$ are continuous on an interval $I$ containing $a$,
  then the IVP has a unique solution on $I$.
\end{theorem}

\begin{example}[Solutions by Inspection]
  For $y''-4y=0$, the functions $e^{2x}$ and $e^{-2x}$ are solutions.
  Every solution has the form
  \[
    y(x)=C_1e^{2x}+C_2e^{-2x}.
  \]
\end{example}

\begin{definition}{Independent Solutions}{independent-solutions}
  Two nonzero solutions $y_1$ and $y_2$ are \textbf{independent} if neither
  is a constant multiple of the other.
\end{definition}

\begin{theorem}{Superposition Theorem}{superposition}
  Suppose $p$ and $q$ are continuous on an interval $I$, and $y_1$, $y_2$
  are independent solutions (Definition~\ref{def:independent-solutions}) of
  \[
    y''+p(x)y'+q(x)y=0.
  \]
  Then all solutions on $I$ have the form
  \[
    y(x)=C_1y_1(x)+C_2y_2(x),\qquad C_1,C_2\in\R.
  \]
\end{theorem}
